% Ported from the chapter-review draft: all VERIFY/COMPLETE markers resolved
% against checkpoint configs and dataset headers (see git history).
\section{Methods}

\subsection{Datasets and study design}

Self-supervised pretraining used 169,120 public methylation profiles represented over 49,156
CpG sites. Profiles were partitioned once into training, validation, and held-out portions
(80/10/10) by a fixed random permutation (seed 42); the held-out portion of 16,912 profiles
was never seen during pretraining and provides the evaluation set for the reconstruction
analyses reported below. The downstream age corpus was derived from the AltuMAge collection
\citep{deLimaCamillo2022} and comprised 10,988 profiles over 21,368 CpG sites, spanning 37
tissue categories and 85 source studies. After removing 328 profiles with implausible age
annotations ($<0$ or $>120$ years) and 71 duplicate profiles, 10,589 remained: 7,177 training,
1,263 validation, and a common held-out test set of 2,149 unique samples. Because the same
test samples were evaluated by all five fold-trained models, variation across models is
reported as training-fold stability rather than as an independent population confidence
interval.

\paragraph{Split integrity.}
Partitions were defined at the level of individual profiles, indexed by their public sample
accession; donor-level grouping was not available for the majority of source studies, which we
note as a limitation. Duplicate and technical-replicate profiles were identified by
methylation fingerprinting across all partitions, yielding 77 duplicate pairs spanning 136
profiles; within each connected component of the resulting graph the lexicographically first
accession was retained and the remaining 71 profiles were excluded. The same exclusion list
and age filter were applied identically to every model and baseline evaluated here, and set
equality of the resulting 2,149 test accessions was verified programmatically rather than
inferred from matching counts. Exact split manifests and their hashes will accompany the
release.

\paragraph{Preprocessing.}
Beta values were restricted to $[0,1]$. Pretraining used the full 49,156-CpG reference panel;
downstream modelling used the 21,368 CpGs measured in the AltuMAge collection. Missing
measurements were excluded from view sampling and from all loss terms through an explicit
validity mask, so that no imputed value ever contributed to training or evaluation. CpGs were
ordered by chromosome and genomic coordinate on the hg38 assembly using the Illumina HM450
manifest. No target-age labels were used during self-supervised pretraining: although the
implementation exposes an optional auxiliary age-regression term, its weight was set to zero
for the pretraining run reported here.

\subsection{Input representation and encoder}

Each observed CpG was represented by its site identity $c_i$ and beta value $b_i$. The token
embedding was
\begin{equation}
  \bm{x}_i = s_c\,\bm{e}_{c_i} + f_{\beta}(b_i),
\end{equation}
where $\bm{e}_{c_i}$ is a learned CpG embedding, $s_c$ is a learned or initialized scale, and
$f_{\beta}$ is a continuous-value encoder based on sinusoidal basis functions followed by a
linear projection. A CLS token was prepended to every partial profile.

The encoder used six transformer layers with hidden size 256, four attention heads (head
dimension 64), feed-forward width 512, and a vocabulary of 49,161 entries (49,156 CpG sites
plus five special tokens). Transformer blocks used
pre-normalization with RMSNorm \citep{Zhang2019}, multi-head self-attention
\citep{Vaswani2017}, SwiGLU feed-forward layers, and rotary position embeddings
\citep{Su2024,Touvron2023}. RoPE position identifiers were assigned using genomic rank rather
than arbitrary input order. The pooled CLS output served as the sample representation.

\subsection{Partial-view WCED pretraining}

WCED adapts the Whole-Cell Expression Decoder introduced for transcriptomic foundation models
\citep{Dandala2025} to array-derived DNA methylation. The original objective reconstructs a
whole-cell expression profile through a CLS bottleneck; our formulation reconstructs CpG
methylation values and evaluates reconstruction only at sites excluded from the encoder input.

The paired-view design was inspired by scConcept, which contrasts representations of different
gene-panel views of the same cell \citep{Bahrami2025}. For each methylation profile, two
independently sampled views $V_i^{(1)}$ and $V_i^{(2)}$ contained 50\% of valid CpGs. A shared
encoder $E_\theta$ produced sample representations
$\bm{h}_i^{(v)}=E_\theta(V_i^{(v)})_{\mathrm{CLS}}$. A shallow decoder $D_\phi$ mapped each
CLS representation to a full-vector prediction $\hat{\bm{b}}_i^{(v)}$.

Reconstruction error was evaluated only over valid CpGs absent from the corresponding input
view, denoted $H_i^{(v)}$:
\begin{equation}
 \mathcal{L}_{\mathrm{recon}} =
 \frac{1}{2}\sum_{v=1}^{2}\frac{1}{|H_i^{(v)}|}
 \sum_{j\in H_i^{(v)}}\left(\tilde b_{ij}-\widetilde{\hat b}_{ij}^{(v)}\right)^2,
\end{equation}
where tildes denote per-sample standardisation: within each view, predicted and observed
values at the withheld positions are centred and scaled to unit variance before the squared
error is taken. This prevents the decoder from reducing the loss by predicting each CpG's
marginal mean, and it makes the reported reconstruction loss scale-free.

A two-layer projection head $g$ ($256\rightarrow128\rightarrow128$, ReLU) mapped CLS
embeddings to $L_2$-normalized vectors $\bm{z}_i^{(v)}$. The symmetric cross-view InfoNCE loss
used matched views as positives and other samples in the batch as negatives \citep{Oord2018}:
\begin{equation}
 \mathcal{L}_{\mathrm{ctr}} = -\frac{1}{2N}\sum_{i=1}^{N}
 \left[
 \log\frac{\exp(\bm{z}_i^{(1)\top}\bm{z}_i^{(2)}/\tau)}
 {\sum_j\exp(\bm{z}_i^{(1)\top}\bm{z}_j^{(2)}/\tau)}+
 \log\frac{\exp(\bm{z}_i^{(2)\top}\bm{z}_i^{(1)}/\tau)}
 {\sum_j\exp(\bm{z}_i^{(2)\top}\bm{z}_j^{(1)}/\tau)}
 \right].
\end{equation}
The complete objective was
$\mathcal{L}=\mathcal{L}_{\mathrm{recon}}+\lambda\mathcal{L}_{\mathrm{ctr}}$, with a fixed
temperature $\tau=0.1$ and $\lambda=0.05$.

\paragraph{Relationship to contrastive single-cell pretraining.}
The contrastive component follows the cell-view formulation of scConcept
\citep{Bahrami2025}, and we state the differences explicitly because they bound
interpretation. First, and most importantly, contrastive alignment is the \emph{sole}
pretraining objective in that framework, which dispenses with reconstruction altogether;
here it is an auxiliary term ($\lambda=0.05$) accompanying withheld-CpG reconstruction. We
adopt the underlying diagnosis---that feature-level reconstruction does not by itself
optimise the sample-level embedding---but address it with two complementary mechanisms:
routing reconstruction through the CLS bottleneck and scoring it only at withheld sites, and
adding explicit cross-view alignment pressure on the same representation. Second, their two
views are disjoint gene panels, whereas our two 50\% CpG subsets are drawn independently and
may overlap, which makes view alignment a comparatively easier task. Third, their negative
set additionally contains embeddings from the anchor's own view, which discourages the
encoder from representing panel identity rather than sample identity; our negative set
contains only cross-view embeddings. Fourth, their temperature is learned during training
whereas ours is fixed, and their contrastive loss is computed on the CLS embedding directly
rather than through a projection head. Finally, their mini-batches are drawn from a single
source dataset so that discrimination cannot be solved using inter-study technical
differences, whereas ours are sampled at random from the pretraining corpus; we return to
this point in the Discussion.

\subsection{Age fine-tuning and comparators}

A two-layer regression head ($256\rightarrow128\rightarrow1$, dropout 0.1) was applied to the
CLS representation. The encoder was frozen for the first ten epochs while the head adapted,
and then unfrozen and optimised jointly at a reduced encoder learning rate
($2\times10^{-5}$ versus $10^{-4}$ for the head). Fine-tuning used the complete measured
profile as input rather than a partial view, and was optimised with a Huber loss on
standardised age. Model selection used validation MedAE without access to the held-out test
labels.

MethylGPT \citep{Ying2024} was evaluated using the same source dataset, age labels, five fold
assignments, and 2,149 held-out sample identifiers. Five predictions per sample were averaged
within each model family before comparison. The pipelines differed in feature handling:
19,608 of the 21,368 downstream CpGs occurred in the common pretraining reference, whereas
the remaining 1,760 CpGs were retained by \method{} under a shared unknown-site identity and
excluded by the MethylGPT conversion pipeline.

As a non-neural reference, we fitted ElasticNetCV \citep{Zou2005} to all 21,368 beta-value
features after standardization based only on the pooled training and validation profiles.
The documented $l_1$-ratio grid (0.1, 0.5, 0.7, 0.9, 0.95, 0.99, and 1.0), an automatically
derived 100-value alpha path, and five-fold internal cross-validation were used once without
post-hoc expansion of the search space. The selected model had $\alpha=0.10883$,
$l_1$ ratio 0.7, and 1,036 non-zero coefficients. The common test set was not used for
hyperparameter selection. A single-fold random-initialization experiment was retained as
exploratory and was not used for the primary comparison; matched objective ablations remain
necessary to attribute downstream performance specifically to WCED.

\subsection{Evaluation and statistical analysis}

Primary endpoints were median absolute error (MedAE) and mean absolute error (MAE), measured
in years. $R^2$ was secondary. Fold-level means and sample standard deviations quantify
sensitivity to the training fold and are not population confidence intervals. For the
\method{}--MethylGPT comparison, 95\% confidence intervals were obtained from 10,000 paired
percentile-bootstrap replicates with random seed 0. Held-out samples, rather than fold-level
predictions, were resampled; the same bootstrap indices were applied to both model ensembles.
The ElasticNet result is reported as a point-estimate reference because per-sample predictions
were not retained for a paired analysis.

Representation analyses included matched-versus-unmatched view cosine similarity, cross-view
retrieval, linear and nonlinear age probes, effective rank, and CLS-versus-mean pooling.
All probe splits were isolated from probe fitting and hyperparameter selection.

\paragraph{Withheld-CpG reconstruction.}
Reconstruction was evaluated on 5,000 profiles drawn from the held-out portion of the
pretraining corpus, which was excluded from pretraining by the fixed partition described
above. Each profile was encoded from a 50\% subset of its measured CpGs, and predictions were
scored only at measured CpGs absent from that input---the same mask used by the training
objective---giving $1.12\times10^{8}$ evaluated positions. Three controls were computed on the
identical profiles under the identical masks: a per-CpG mean predictor, the decoder applied to
CLS embeddings shuffled across the batch, and the decoder applied to standard normal vectors.
The per-CpG means were computed on the evaluation cohort itself, which favours that baseline
and therefore makes the comparison conservative.

\paragraph{Cross-view consistency.}
Two 50\% views were encoded for each of 2,000 profiles and compared by cosine similarity of the
resulting CLS embeddings. We report the mean similarity for matched and unmatched pairs, and
retrieval accuracy at rank $k$: the fraction of first views whose matched second view falls
within the $k$ nearest of all 2,000 candidates. Retrieval at $k=1$ is the strictest point on
this curve; chance is $k/2000$. The whole evaluation was repeated for five independent random
view draws, and we report the mean and range across draws rather than a single draw.

Two view constructions were compared. In the first, matching the construction used during
pretraining, the two views are drawn independently and therefore share approximately half of
their CpGs in expectation. In the second, they form a random partition of the valid CpGs and
share no site, following the disjoint-panel construction of \citet{Bahrami2025}. The disjoint
condition is out of distribution for an encoder trained under the overlapping construction, and
we therefore interpret it as a conservative bound on cross-subset consistency rather than as a
matched comparison. Both conditions used the same profiles, checkpoint, and encoder call, so
they differ only in how the two views are formed.
